\section{KẾT LUẬN}

Đề tài đã xây dựng được một prototype Self-Evolving AI Agents cho xử lý sự cố mạng viễn thông trên nền tảng NIKA. Hệ thống tích hợp môi trường Kathará, các MCP server chẩn đoán, workflow ReAct/agent và hai cơ chế tự tiến hóa gồm Procedural Memory và Tool Refinement. Kết quả thực nghiệm trên hai lần chạy 125 case chung cho thấy hướng tiếp cận tự tiến hóa giúp tăng incident success rate từ 20.8\% lên 28.0\%, tăng RCA F1 từ 0.304 lên 0.387, đồng thời giảm số bước và số lần gọi công cụ trung bình.

Tính mới của đề tài nằm ở việc không xem agent như một thành phần tĩnh, mà thiết kế vòng lặp học sau mỗi episode. Tính sáng tạo thể hiện ở sự kết hợp giữa bộ nhớ thủ tục có điều kiện kích hoạt, truy xuất theo thuộc tính miền mạng và tinh chỉnh tài liệu công cụ dựa trên trace thực tế. Tính hiệu quả được minh chứng bằng cải thiện định lượng ở các metric chẩn đoán quan trọng, đặc biệt là RCA và tỉ lệ giải quyết incident hoàn chỉnh.

Về đóng góp cá nhân, sinh viên đã trực tiếp nghiên cứu các hướng self-evolving agent, thiết kế kiến trúc tích hợp với NIKA, hiện thực hai mô-đun học kinh nghiệm, xây dựng pipeline benchmark và phân tích kết quả. Điểm quan trọng của sản phẩm là các thành phần tiến hóa đều có trạng thái rõ ràng, có thể kiểm tra lại và có thể tắt/bật để đánh giá ablation trong các thí nghiệm tiếp theo.

Trong thực tế, hệ thống có thể được phát triển thành trợ lý cho kỹ sư vận hành mạng: tự thu thập bằng chứng, đề xuất nguyên nhân, giải thích quá trình suy luận và học từ phản hồi chuyên gia. Các hướng phát triển tiếp theo gồm tăng khả năng tổng quát hóa sang topology/giao thức mới, bổ sung cơ chế human-in-the-loop để phê duyệt kỹ năng trước khi lưu lâu dài, tối ưu token khi nạp bộ nhớ và mở rộng benchmark sang multi-fault incidents có nhiều nguyên nhân đồng thời. Một hướng quan trọng khác là xây dựng dashboard giám sát quá trình tiến hóa, cho phép xem kỹ năng nào được kích hoạt, tài liệu công cụ nào đã hội tụ và episode nào đóng góp vào thay đổi bộ nhớ.
